% Dialect Classification Evidence — Thematic Synthesis
% !TEX program = xelatex
% !TEX encoding = UTF-8 Unicode

\section{Dialect Classification Evidence}
\label{sec:dialect}

The SNEA languages belong to different dialect groups within the Eastern
Algonquian family. The corpus data provides evidence for these
classifications, and understanding the dialect relationships is essential
for cross-source calibration. A form that looks like a transcription error
in one dialect may be a regular phonological reflex in another.

\subsection{The Four-Way Dialect Classification}

Southern New England Algonquian languages are traditionally classified by
the reflex of Proto-Algonquian *r (sometimes reconstructed as *l or *θ).
Four dialect types are recognized:

\begin{table}[h]
\centering
\caption{SNEA dialect classification by PA *r reflex}
\label{tab:dialect-classification}
\begin{tabular}{llll}
\toprule
\textbf{Dialect type} & \textbf{Reflex of PA *r} & \textbf{Languages} & \textbf{Sources} \\
\midrule
\textbf{N-dialect} & /n/ & Massachusett (Wampanoag), & Eliot, Trumbull, \\
 & & Narragansett, Nipmuc & Williams \\
\textbf{Y-dialect} & /j/ & Mohegan-Pequot, Montauk & Prince-Speck, Fielding \\
\textbf{R-dialect} & /r/ & Western Connecticut (Loup A) & (limited) \\
\textbf{L-dialect} & /l/ & Munsee Delaware, Unami & Zeisberger \\
\bottomrule
\end{tabular}
\end{table}

The PA *r reflex is the most diagnostic isogloss for SNEA dialect
classification, but it is not the only one. Costa (2007) identifies a
network of isoglosses that collectively define the SNEA dialect continuum.

\subsection{Costa (2007) Isoglosses}

David Costa's definitive study ``The Dialectology of Southern New England
Algonquian'' (2007) identifies the following key isoglosses that distinguish
the SNEA languages:

\begin{enumerate}
  \item \textbf{PA *r reflex:} /n/ (N-dialect) vs.\ /j/ (Y-dialect) vs.\
    /r/ (R-dialect) vs.\ /l/ (L-dialect). This is the primary
    classification criterion.

  \item \textbf{PA *θ reflex:} In N-dialects, PA *θ merges with /n/ (e.g.,
    Massachusett \textit{neen} 'I' from PA *niːra, where *r $\rightarrow$ /n/
    and *θ $\rightarrow$ /n/). In Y-dialects, PA *θ $\rightarrow$ /j/ in some
    environments.

  \item \textbf{PA *š reflex:} In Mahican, PA *š merges with /s/ (a
    distinguishing feature from the other SNEA languages, where /š/ is
    retained as a separate phoneme).

  \item \textbf{Word-final vowel deletion:} Mahican shows more advanced
    word-final vowel deletion than the other SNEA languages, reflecting its
    transitional status toward Delaware.

  \item \textbf{Preaspiration retention:} Mohegan-Pequot (as recorded by
    Prince-Speck) appears to have retained preaspiration longer than
    Narragansett or Massachusett, though this may be a recording artifact.

  \item \textbf{Stress patterns:} Narragansett stress patterns (as recorded
    by Williams' diacritics) appear similar to Loup A and Massachusett,
    with a preference for stress on long syllables and alternating rhythm on
    short syllables.
\end{enumerate}

\subsection{Mohegan-Pequot as a Y-Dialect: The Evidence}

Mohegan-Pequot is classified as a Y-dialect based on the reflex of PA *r
$\rightarrow$ /j/. The evidence from the corpus includes:

\begin{itemize}
  \item \textbf{PA *r $\rightarrow$ /j/ in cognate sets:} Where Massachusett
    (N-dialect) has /n/, Mohegan-Pequot has /j/. Stephanie Fielding noted
    this correspondence operating even in modern Connecticut English
    pronunciation (e.g., ``Sullivan'' $\rightarrow$ ``Suh-yuh-van'').

  \item \textbf{The Y-dialect isogloss bundle:} Costa (2007) shows that the
    Y-dialect area (Mohegan-Pequot, Montauk) forms a coherent isogloss
    bundle that separates it from the N-dialect area (Massachusett,
    Narragansett, Nipmuc).

  \item \textbf{Prince-Speck evidence:} The Prince-Speck glossary (1904)
    confirms the Y-dialect classification for early 20th-century
    Mohegan-Pequot, showing /j/ reflexes where N-dialects have /n/.

  \item \textbf{Fielding evidence:} The modern Mohegan dictionary (2013)
    confirms the Y-dialect classification for the revived language.
\end{itemize}

Granberry (2003:13) explicitly documents the core correspondences: Pequot
\textit{p, t, k, ch, s} $\rightarrow$ Mohegan \textit{b, d, g, j, z}

\subsection{Central Algonquian Reference Points}

The dialect classification can be further verified against Central Algonquian
languages, which preserve the PA *r reflex distinctions more clearly than
the SNEA colonial sources:

\begin{itemize}
  \item \textbf{Cree (Y-dialect):} Plains Cree (\url{https://itwewina.altlab.app})
    is a Y-dialect (PA *r $\rightarrow$ /y/), providing a reference for the
    Mohegan-Pequot Y-dialect classification. Cree also preserves the PA
    vowel system more faithfully than any SNEA source.
  \item \textbf{Ojibwe (N-dialect):} The Ojibwe People's Dictionary
    (\url{https://ojibwe.lib.umn.edu}) represents an N-dialect (PA *r
    $\rightarrow$ /n/), providing a reference for the Wampanoag and
    Narragansett N-dialect classification.
  \item \textbf{Munsee Delaware (L-dialect):} The Lenape Talking Dictionary
    (\url{https://www.talk-lenape.org}) represents an L-dialect (PA *r
    $\rightarrow$ /l/), providing a reference for the Mahican transitional
    status between SNEA and Delaware.
\end{itemize}

These Central Algonquian reference points confirm that the SNEA dialect
classifications are consistent with the broader Algonquian family patterns.
The Y/N/L/R classification is not unique to SNEA --- it reflects a
continent-wide pattern of PA *r reflexes across the Algonquian family.
form-initially, and final Pequot \textit{p, t, k, ch, s} $\rightarrow$
Mohegan form-final \textit{b, d, g, j, z}. This lenis-fortis alternation is
a distinctive feature of the Mohegan-Pequot dialect complex.

\subsection{Mahican as Transitional}

Mahican occupies a transitional position between the SNEA languages and the
Delaware languages. Evidence for this transitional status includes:

\begin{itemize}
  \item \textbf{PA *š $\rightarrow$ /s/:} Mahican merges PA *š with /s/, a
    feature shared with Delaware but not with the other SNEA languages
    (which retain /š/ as a separate phoneme).

  \item \textbf{PA *θ $\rightarrow$ /n/:} Mahican merges PA *θ with /n/,
    like the N-dialects of SNEA, but unlike Delaware (which retains /θ/ or
    shifts it to /s/).

  \item \textbf{Word-final vowel deletion:} Mahican shows more advanced
    word-final vowel deletion than the other SNEA languages, approaching the
    Delaware pattern.

  \item \textbf{PA *eː $\rightarrow$ /a/:} Mahican retracts/lowers PA *eː to
    /a/, a development shared with some Delaware dialects but not with the
    other SNEA languages (where PA *eː $\rightarrow$ /iː/).

  \item \textbf{Edwards' data:} Edwards' Mahican forms show features of
    both groups. His \textit{neesoh} 'two' shows the SNEA /n/ reflex of PA
    *r, but the final \textit{-oh} suffix is more characteristic of Delaware.
\end{itemize}

This transitional status means that Mahican cannot be used as a direct model
for the other SNEA languages without accounting for the differences. The
conversion workflow must apply different sound correspondences for Mahican
than for the N-dialect or Y-dialect languages.

\subsection{Cross-Linguistic Cognate Evidence Table}

Table~\ref{tab:cognate-dialect} shows how the same PA root is reflected
across the four dialect types, illustrating the regular sound
correspondences that define the classification.

\begin{table}[h]
\centering
\caption{Cross-dialect cognate evidence for PA *r reflex}
\label{tab:cognate-dialect}
\small
\begin{tabular}{llllll}
\toprule
\textbf{Gloss} & \textbf{PA} & \textbf{N-dialect} & \textbf{Y-dialect} & \textbf{Mahican} & \textbf{L-dialect} \\
 & & \textbf{(Mass.)} & \textbf{(Mohegan)} & \textbf{(Edwards)} & \textbf{(Munsee)} \\
\midrule
'I' & *niːra & neen & neen & neen & niːl \\
'thou' & *kiːra & keen & keen & keen & kiːl \\
'man' & *naːpeːwa & napp & nap & --- & náːpeːw \\
'house' & *wiːkiwaːmi & wetu & wetu & --- & wiːkiwaːm \\
'water' & *nepyi & nippe & nip & --- & mbíː \\
'dog' & *aɬemwa & arum & ayum & --- & --- \\
'three' & *neʔθwi & shwa & shwa & --- & nəxá \\
'four' & *nyeːwi & yaw & yaw & --- & néːwi \\
'five' & *nyaːlanwi & pàpàna & --- & nannah & náːlan \\
\midrule
\textbf{PA *r reflex} & & \textbf{/n/} & \textbf{/j/} & \textbf{/n/} & \textbf{/l/} \\
\bottomrule
\end{tabular}
\end{table}

The table reveals several important patterns:

\begin{enumerate}
  \item \textbf{The PA *r reflex is the primary classifier.} Massachusett
    (N-dialect) and Mohegan (Y-dialect) both show /n/ in 'I' and 'thou',
    but the underlying PA source differs: Massachusett /n/ comes from PA *r,
    while Mohegan /n/ comes from a different source (PA *n). The Y-dialect
    reflex of PA *r is /j/, which appears in other lexical items.

  \item \textbf{Mahican patterns with N-dialects for *r.} Edwards' Mahican
    shows /n/ for PA *r, aligning it with the N-dialect group for this
    isogloss.

  \item \textbf{Munsee (L-dialect) is clearly distinct.} The L-dialect
    reflex /l/ for PA *r separates Munsee from all the SNEA languages.

  \item \textbf{The PA *ɬ reflex distinguishes Narragansett from
    Massachusett.} The 'dog' cognate shows Narragansett /j/ (\textit{ayim})
    vs.\ Massachusett /r/ (\textit{arum}), both from PA *aɬemwa. This is a
    secondary isogloss within the N-dialect group.

  \item \textbf{Lexical stability across dialect boundaries.} Basic
    vocabulary items like 'house', 'water', and 'three' are remarkably
    similar across all four dialect types, confirming the reliability of
    these items for cross-source calibration.
\end{enumerate}

\subsection{Implications for the Conversion Workflow}

The dialect classification has direct implications for the conversion
workflow:

\begin{enumerate}
  \item \textbf{Source-specific sound correspondences.} Each source language
    requires its own set of PA $\rightarrow$ daughter sound correspondences.
    Applying Massachusett correspondences to Mohegan-Pequot data would
    produce systematic errors.

  \item \textbf{Cross-dialect calibration is valid only for shared
    isoglosses.} When using Massachusett to disambiguate Narragansett (both
    N-dialects), the calibration is valid for shared features. When using
    Massachusett to disambiguate Mohegan-Pequot (N-dialect vs.\ Y-dialect),
    the calibration must account for the isogloss differences.

  \item \textbf{Mahican requires separate treatment.} Mahican's transitional
    status means it cannot be grouped with either the N-dialect or Y-dialect
    languages for conversion purposes. A separate set of sound
    correspondences is required.

  \item \textbf{The cognate set approach is the most reliable calibration
    method.} By anchoring each conversion in a known PA reconstruction and
    applying the correct dialect-specific sound correspondences, the
    conversion workflow can produce reliable phonemic transcriptions even
    from highly ambiguous colonial orthographies.
\end{enumerate}

\subsection{Cross-Linguistic Reassessment of Dialect Classification}

The newly acquired reference dictionaries provide independent evidence that
both confirms and complicates the traditional four-way dialect classification:

\begin{itemize}
  \item \textbf{Watkins Cree dialect chain model:} Watkins (1865) documents
    a l/n/th/y dialect chain across Cree, where the same PA segment surfaces
    as /l/, /n/, /θ/, or /j/ depending on dialect. This is a continuum, not a
    discrete classification. The SNEA Y/N/R/L classification may similarly
    represent a continuum rather than discrete groups --- the apparent
    binary between Mohegan (Y-dialect) and Wampanoag (N-dialect) may be an
    artifact of limited documentation rather than a genuine linguistic
    boundary. The Cree model suggests that intermediate dialects (R-dialect
    Loup A, L-dialect Munsee) are not separate types but points on a
    continuum.

  \item \textbf{Brinton Lenape Unami/Minsi split:} Brinton (1888) documents
    a dialect split within Lenape (Unami vs.\ Minsi) that parallels the
    SNEA N-dialect/Y-dialect split. The Unami/Minsi distinction involves the
    same PA *r reflex (Unami /l/ vs.\ Minsi /n/), confirming that the
    Y/N/R/L classification is a recurring pattern across Eastern Algonquian,
    not unique to SNEA.

  \item \textbf{Sauk vowel system as PA model:} The Sauk 4-vowel length
    system (Whittaker 2005) provides a template for the PA vowel system that
    underlies all dialect classifications. The fact that Sauk (Central
    Algonquian) preserves a clean 4-way length distinction suggests that
    the SNEA languages' original vowel systems were similarly structured,
    and that the ambiguity in colonial transcriptions reflects recording
    limitations rather than phonological merger.

  \item \textbf{Arapaho divergence as boundary condition:} The Arapaho
    Dictionary (Cowell et al.\ 2012) represents the extreme end of Algonquian
    phonological divergence. Arapaho's radical changes (PA *p $\rightarrow$
    /b/, PA *k $\rightarrow$ /g/ or /k/, development of /θ/, /x/, triple
    vowels) provide the outer bounds for what is possible in Algonquian sound
    change. The SNEA languages, by contrast, are relatively conservative ---
    their sound changes are minor compared to Arapaho's transformations.
\end{itemize}

These comparative data suggest that the traditional four-way classification
is essentially correct but should be understood as a continuum rather than
discrete categories. The conversion workflows should treat the dialect
classification as a gradient feature, with each source occupying a position
on the continuum rather than belonging to a discrete type.
